\subsection{定义}\label{sec:9j}
为定义带角正合辛流形 $M$ 的 Fukaya 范畴初步版本 $\F(M)^{\mathrm{pr}}$，需要选取以下额外数据：
\begin{itemize}
\item 一致的条带状端万有选择 $\{\epsilon_k^{d+1}\}$；
\item 对 $M$ 中每一对正合 Lagrangian 子流形 $(L_0,L_1)$ 选定 Floer 数据 $(H_{L_0,L_1},J_{L_0,L_1})$；
\item 一致的扰动数据万有选择 $(K_{L_0,\ldots,L_d},J_{L_0,\ldots,L_d})$，与先前选择的条带状端及 Floer 数据相容。
\end{itemize}
这些数据还须满足一个正则性条件，现在来说明。对于正合 Lagrangian 子流形的每个 $(d+1)$ 元组 $(L_0,\ldots,L_d)$（$d\geq2$），以及点 $y_0\in\mathcal C(L_0,L_d)$、$y_k\in\mathcal C(L_{k-1},L_k)$（$1\leq k\leq d$），考虑
\begin{equation}\label{eq:9.22}
\M^{d+1}(y_0,\ldots,y_d)=\M_{\mathcal S^{d+1}}(y_0,\ldots,y_d).
\end{equation}
这是族 $\mathcal S^{d+1}\to\mathcal R^{d+1}$ 上方程\eref{8.9}的解所成的模空间\eref{9.17}；该族带有 Lagrangian 标记 $(L_0,\ldots,L_d)$ 和已选定的扰动数据，渐近条件是在条带状端上收敛于 $y_k$。
\begin{definition*}
称\eref{9.22}中的点 $(r\in\mathcal R^{d+1},u:\mathcal S_r^{d+1}\to M)$ 为\emph{非齐次伪全纯多边形}。将记号延伸到 $d=1$，令
\SourcePage{135}
$\M^2(y_0,y_1)=\M_Z^*(y_0,y_1)$，即 Floer 方程的解模去平移后的模空间。
\end{definition*}
所要求的条件是：这些模空间全部正则。对于 $d=1$，已经提到，这对一般位置的 Floer 数据选择成立。对于 $d\geq2$，在一致的扰动数据万有选择这一类别中，同样的结论成立。
\begin{definition*}
$\F(M)^{\mathrm{pr}}$ 的对象是正合 Lagrangian 子流形 $L\subset M$，态射由 Floer 上链空间给出：
\[
\hom_{\F(M)^{\mathrm{pr}}}(L_0,L_1)=CF^{\mathrm{pr}}(L_0,L_1).
\]
一阶映射 $\mu^1=\partial$ 是边界算子，像\eref{8.19}一样，由计数 $\M^2(y_0,y_1)$ 中的孤立点定义。高阶复合类似地定义为
\begin{equation}\label{eq:9.23}
\mu^d(y_d,\ldots,y_1)=\sum_{y_0}\#\M^{d+1}(y_0,\ldots,y_d)\,y_0.
\end{equation}

\end{definition*}
\znote{127}{这个定义只覆盖了“横截”的情形：态射空间 $CF^{\mathrm{pr}}(L_0,L_1)$ 由交点生成，而交点本身依赖于所选的 Floer 数据，也就是依赖于对 $L_0,L_1$ 所作的 Hamilton 扰动。严格说来，$\mu^d$ 不是 $(L_0,\ldots,L_d)$ 的函数，而是后者连同所选数据的函数，§(10a) 用总范畴与协调系统来处理这种依赖。当两个 Lagrangian 子流形不横截相交时，交点不是离散集，计数失去意义；一般情形下 Fukaya 范畴的定义之所以困难，症结正在这里。[来源：nLab "Fukaya category"，https:/\allowbreak /\allowbreak ncatlab.org/\allowbreak nlab/\allowbreak show/\allowbreak Fukaya+\allowbreak category（2026-11-15 访问）。]}
紧致性及粘合论证表明，这一定义良好：每个 $\M^{d+1}(y_0,\ldots,y_d)$ 的零维部分都是有限集；并且满足 $\Ai$ 结合律方程\eref{1.2}。\znote{127}{这里的 $\#$ 指零维模空间的点数按系数域 $\K$ 取的模 $2$ 计数（\stref{笔记}{9.7}始终假定 $\operatorname{char}\K=2$），因此\eref{9.23}的右端是有限和。在分次版本中，$\mu^d$ 的次数为 $2-d$，于是\eref{1.2}中每一项的总次数一致；这与 $\Ai$ 代数通常的约定相同，即 $D_k$ 的次数为 $2-k$（$D_1$ 是微分、$D_2$ 是乘法、$D_3$ 是结合子）。[来源：nLab "A-∞ algebra"，https:/\allowbreak /\allowbreak ncatlab.org/\allowbreak nlab/\allowbreak show/\allowbreak A-infinity-algebra（2026-11-15 访问）。]}因此，$\F(M)^{\mathrm{pr}}$ 是一个不分次的 $\Ai$ 范畴。按定义，$\mu^2$ 在上同调上诱导的映射是三角形乘积\eref{8.3}，所以底层的上同调范畴确实就是先前定义的 $H\F(M)^{\mathrm{pr}}$。这也说明 $\F(M)^{\mathrm{pr}}$ 上同调含幺。\znote{127}{“上同调含幺”指 $H^0(\hom(L,L))$ 中存在元素 $e_L$，使 $\mu^2(e_L,a)=a$、$\mu^2(a,e_L)=a$ 在上同调层成立。它弱于链层的含幺性（存在闭元素 $e_L$ 使 $\mu^2(e_L,-)=\mu^2(-,e_L)=\id$ 在链层成立），因为 $\mu^2$ 只在上同调层才结合，链层上结合律只能到同伦。这里的 $e_L$ 来自 $L$ 与它的小 Hamilton 扰动像的横截交点中最靠近对角线的那一点。[来源：原书 §(1a) 关于上同调范畴与“非含幺”术语的约定，见第一、三处定义。]}

\begin{note}\label{st:9.7}
我们没有明确地定义不分次 $\Ai$ 范畴，但其含义应当清楚。回顾，假设系数域满足 $\operatorname{char}(\K)=2$，所以定义符号时不需要分次。有了这一点，\sr{chapter:I}{第一章}的初等理论只需作表面上的调整便仍然适用，唯一的例外是\stref{推论}{5.8}。读者应当记住，目前的设置只是暂时的简化：最终所有应用都会在分次有良好定义的背景下进行。
\end{note}

\begin{note}\label{st:9.8}
上一段概述的结果为本领域专家所熟知，证明虽然有些长，但并未超出伪全纯曲线理论的标准技术。在本节余下部分，将简要勾勒主要步骤，希望由此传达这些技术的大致面貌。对希望深入细节的读者，以下是一些可能的参考指引。Riemann 曲面族的非齐次项似乎最早用于\bcite{116}。横截性论证的基本参考是\bcite{49}。Floer 的原始论文\bcite{46}、\bcite{45}、\bcite{44}，以及\bcite{103}、\bcite{50}，包含处理 Lagrangian 边界条件所需的一些基本分析。空间 $\mathcal R^{d+1}$ 及其紧化在\bcite{58}中讨论，该文还包含许多与 Fukaya 范畴有关的材料。整套构造中最复杂的部分，
\SourcePage{136}
可能是用来证明 $\Ai$ 结合律方程成立的粘合过程；见\bcite{60}，第 18 节。这个论证推广了 Floer 理论中更熟悉的 $\partial^2=0$ 的证明；其 Lagrangian 版本和 Hamilton 版本分别见\bcite{44}和\bcite{119}。
\end{note}

\subsection{横截性}\label{sec:9k}
固定一致的条带状端万有选择，并为每一对正合 Lagrangian 子流形固定 Floer 数据。假设已经为所有 $d\leq d_0$ 选取扰动数据\eref{9.20}。取一组 Lagrangian 标记 $(L_0,\ldots,L_{d_0})$。希望稍作修改给定扰动数据 $(K,J)=(K_{L_0,\ldots,L_{d_0}},J_{L_0,\ldots,L_{d_0}})$，使对应模空间\eref{9.22}正则。形式上，$(K,J)$ 的无穷小变形由一对
\begin{equation}\label{eq:9.24}
(\delta K,\delta J)\in\Omega^1_{\mathcal S^{d_0+1}/\mathcal R^{d_0+1}}(\mathcal S^{d_0+1},\mathcal H)
\times C^\infty(\mathcal S^{d_0+1},T_J\mathcal J)
\end{equation}
给出，要求 $(\delta K,\delta J)$ 在每个曲面 $\mathcal S_r^{d_0+1}$ 的条带状端上为零；并且对连通分支 $C_k\subset\partial\mathcal S_r^{d_0+1}$ 的每个切向量 $\xi$，有 $\delta K(\xi)|_{L_k}=0$。这样的无穷小变形很容易“取指数”成为实际的变形：
\begin{equation}\label{eq:9.25}
(K',J')=(K+\delta K,\ J\exp(-J\cdot\delta J)).
\end{equation}
为了保持与 $d<d_0$ 时的扰动数据选择的一致性，需要对容许的变形空间作一些限制。对每棵稳定 $d_0$ 叶树 $T$，固定粘合参数充分小的开子集 $U^T\subset(-1,0)^{\operatorname{Ed}^{\mathrm{int}}(T)}\times\mathcal R^T$。只要各 $U^T$ 取得充分小，就可选开子集 $\Omega\subset\mathcal S^{d_0+1}$，使对每个 $r\in\mathcal R^{d_0+1}$，交集 $\Omega_r=\Omega\cap\mathcal S_r^{d_0+1}$ 非空，并具有以下性质：$\Omega_r$ 与 $\mathcal S_r^{d_0+1}$ 的条带状端不交；若对某棵 $T$ 有 $r\in\gamma^T(U^T)$，则 $\Omega_r$ 位于 $\mathcal S_r^{d_0+1}$ 对应分解的厚部。\znote{128}{把 $\Omega_r$ 限制在厚部，是因为薄部与条带状端上的扰动数据已被低阶数据决定：按 §(9i) 的一致性条件，若 $r$ 落在某个 $\gamma^T(U^T)$ 中，$\mathcal S_r^{d_0+1}$ 就带有由 $d<d_0$ 阶数据粘合而成的结构，那里不能自由改动。厚部是圆盘“真正有 $d_0+1$ 个标记”的部分，扰动数据在那里仍可任意选取；\eref{9.27}之后用 $\delta$ 函数逼近的论证正依赖这一点。[来源：原书 §(9i) 的一致性条件与 §(9e) 的厚薄分解，见\eref{9.19}、\eref{9.20}。]}

考虑\eref{9.24}中各对 $(\delta K,\delta J)$ 所成的子空间 $\mathcal T$，要求它们在 $\Omega$ 外为零，并相对于部分紧化 $\overline{\mathcal S}^{d_0+1}$ 渐近衰减。最后一个条件意味着，$(\delta K,\delta J)$ 光滑延伸为 $\overline{\mathcal S}^{d_0+1}$ 上的一对 $(\overline{\delta K},\overline{\delta J})$，它及其任意阶导数都在 $\overline{\mathcal S}^{d_0+1}\setminus\mathcal S^{d_0+1}$ 上为零。对于每个这样的 $(\delta K,\delta J)$，变形后的扰动数据\eref{9.25}仍满足关于粘合映射 $\gamma^T$ 的一致性。具体说，因 $(K',J')=(K,J)$ 在 $\Omega$ 外成立，曲面 $\mathcal S_r^{d_0+1}$ 薄部上的条件仍满足；一致性要求的另一部分涉及相对于 $\overline{\mathcal S}^{d_0+1}$ 的渐近行为，而 $(K',J')$ 与原 $(K,J)$ 的渐近行为也相同。

\begin{proof}[正则性论证]
还需证明，这类变形足以使带有标记 $(L_0,\ldots,L_{d_0})$ 的非齐次伪全纯 $(d_0+1)$ 边形模空间正则。这是 Sard–Smale 定理的标准应用。
\SourcePage{137}
对于原扰动数据 $(K,J)$ 的模空间中任意一点 $(r,u)\in\M^{d_0+1}(y_0,\ldots,y_{d_0})$，都有一个万有线性化算子 $D_{\mathcal T,\mathcal S^{d_0+1},r,u}$。它包含\eref{9.18}，但另有变量，用来表示修改 $(K,J)$ 时 $(du-Y)^{0,1}$ 如何变化。为严格起见，应当为族 $\mathcal S^{d_0+1}$ 选择局部平凡化，见\stref{笔记}{9.4}；为简便起见，这里略去。具体公式为
\begin{equation}\label{eq:9.26}
\begin{aligned}
\overline D_{\mathcal S^{d_0+1},r,u}:&\ \mathcal T\times T_r\mathcal R^{d_0+1}
\times(T\mathcal B_{\mathcal S_r^{d_0+1}})_u\longrightarrow(\mathcal E_{\mathcal S^{d_0+1},r})_u,\\
(\delta K,\delta J,\rho,X)\longmapsto{}&-(\delta Y)^{0,1}
+\tfrac12\delta J\circ(du-Y)\circ I_S+D_{\mathcal S^{d_0+1},r,u}(\rho,X),
\end{aligned}
\end{equation}
其中 $\delta Y$ 是由 $\delta K$ 得到的、取值于 Hamilton 向量场的逐纤维一形式，$I_S$ 是 $S=\mathcal S_r^{d_0+1}$ 上的复结构。主要步骤是证明\eref{9.26}为满射。注意，由于最后一项是 Fredholm 算子，像自动闭合。\eref{9.26}余核中的元素可表示为
\[
\chi\in L^q(S,(\Omega_S^{0,1}\otimes u^*TM)^\vee),\qquad 1/p+1/q=1,
\]
特别地，它必须满足 $D_{S,u}^*\chi=0$。由椭圆正则性，$\chi$ 至少在 $\operatorname{int}(S)$ 中光滑；更仔细考察其性质，实际上可证明它直到边界都光滑，但这里不需要这一点。此外，对所有 $\delta Y$，它满足
\begin{equation}\label{eq:9.27}
\int_S\langle(\delta Y)^{0,1},\chi\rangle=0.
\end{equation}
由构造，存在非空开集 $\Omega_r\subset S$，在其中 $\delta K$ 的选择基本上任意。特别地，对于任意 $z\in\Omega_r\cap\operatorname{int}(S)$ 及任意 $\alpha\in\Omega^{0,1}(S)_z\otimes T M_{u(z)}$，存在一列 $(\delta K,\delta J)$，使相应的 $(\delta Y)^{0,1}$ 弱收敛于 $\delta$ 函数 $\delta_z\otimes\alpha$。将这一列代入\eref{9.27}，可见 $\chi_z=0$。因此，$\chi$ 在 $\Omega_r$ 上为零；它处处为零，则由方程
\[
D_{S,u}^*\chi=0
\]
的解的唯一延拓性质推出。
\end{proof}

\begin{note}\label{st:9.9}
刚才给出的证明概要，至少略去了一个例行却非平凡的技术点。容许的 $(\delta J,\delta K)$ 所成的空间 $\mathcal T$ 是 Fréchet 空间。为了应用 Sard–Smale 定理，需要进一步限制到一个稠密 Banach 子空间，例如采用 Floer 的 $C_\epsilon^\infty$ 拓扑的一种版本，将关于 $\overline{\mathcal S}^{d_0+1}$ 所要求的渐近行为计入其中。这样，\eref{9.26}的定义域在前两个变量上缩减为稠密子空间。然而，因为仍保留第三部分所给出的 Fredholm 算子，满射性不受影响。
\end{note}

\clearpage
\subsection{紧致性}\label{sec:9l}
非齐次伪全纯多边形的模空间有自然的 Gromov 紧化，其结构在一定程度上与 $\mathcal R^{d+1}\subset\overline{\mathcal R}^{d+1}$ 的结构相仿。实际上，后者可以视为
\SourcePage{138}
$M=L_0=\cdots=L_d=\text{一点}$ 的平凡情形。下面先在集合意义下定义这些紧化，不加证明地陈述其主要拓扑性质，再说明这如何推出各映射 $\mu^d$ 构成一个 $\Ai$ 结构。

像\eref{9.22}的定义中那样固定 $L_k,y_k$，再固定一棵半稳定 $d$ 叶树 $T$。假设对每个顶点 $v\in\operatorname{Ve}(T)$ 都给定一个 $|v|$ 标点圆盘 $S_v$。若 $|v|\geq3$，该圆盘确定一点 $r_v\in\mathcal R^{|v|}$，并存在唯一同构 $S_v\cong\mathcal S_{r_v}^{|v|}$。在余下的 $|v|=2$ 情形，有 $S_v\cong Z$，但这一同构只在模去 $s$ 方向平移的意义下确定。此外，要求这些圆盘带有 Lagrangian 标记，使每个旗标 $f=(v,e)$ 都有一对指定子流形 $(L_{f,0},L_{f,1})$。像前面讨论粘合时一样，标记应彼此相容，使粘合后的 $(d+1)$ 标点圆盘上的标记为给定的 $L_k$。最后，对每个 $f$，要求有一点 $\tilde y_f\in\mathcal C(L_{f,0},L_{f,1})$，满足 $\tilde y_{f^+(e)}=\tilde y_{f^-(e)}$，且与根、叶相邻的旗标所对应的点，按自然次序为 $y_0,\ldots,y_d$。
\begin{definition*}
\emph{非齐次伪全纯稳定多边形}，是一组映射 $u_v:S_v\to M$（$v\in\operatorname{Ve}(T)$），满足以下条件：
\begin{itemize}
\item 若 $|v|\geq3$，可将 $u_v$ 看成映射 $\mathcal S_{r_v}^{|v|}\to M$，要求它是对应的非齐次伪全纯映射方程的解，即
\[
(r_v,u_v)\in\M^{|v|}(\tilde y_{f_0(v)},\ldots,\tilde y_{f_{|v|-1}(v)}),
\]
其中 $f_0(v),\ldots,f_{|v|-1}(v)$ 是按自然次序排列的 $v$ 周围的旗标，扰动数据取自我们的万有选择。
\item 若 $|v|=2$，则 $u_v:Z\to M$ 应当代表
\[
u_v\in\M^2(\tilde y_{f_0(v)},\tilde y_{f_1(v)}),
\]
即它是给定 Floer 数据的 Floer 方程的非驻定解。因为该方程具有平移不变性，这一表述不依赖于同构 $S_v\cong Z$ 的选择。
\end{itemize}
\end{definition*}
定义 $\M^T(y_0,\ldots,y_d)$ 为所有这种非齐次伪全纯稳定多边形的等价类的集合。等价关系是各带标点圆盘 $S_v$ 的同构，并要求各映射 $u_v$ 在此同构下对应相等。具体说，若某个 $v$ 满足 $|v|=2$，对相应的 $u_v$ 作平移，并不改变点 $\{u_v\}\in\M^T(y_0,\ldots,y_d)$。在集合意义下，将 Gromov 紧化定义为对所有 $T$ 的不交并：
\begin{equation}\label{eq:9.28}
\overline\M^{d+1}(y_0,\ldots,y_d)=\coprod_T\M^T(y_0,\ldots,y_d).
\end{equation}
与\eref{9.11}不同，半稳定树 $T$ 有无限多棵，因为总能继续插入二价顶点。然而，无论怎样选取 $(y_0,\ldots,y_d)$，\eref{9.28}中实际非空的层
\SourcePage{139}
只有有限多个。为说明这一点，对各分量的能量\eref{8.12}求和，得到
\[
\sum_vE(u_v)=A_{H_{L_0,L_d}}(y_0)-\sum_{k=1}^dA_{H_{L_{k-1},L_k}}(y_k)+(\text{曲率项}).
\]
曲率项只由 $|v|\geq3$ 的顶点贡献；考察扰动数据万有选择的一致性条件，可以推出它在整个空间\eref{9.28}上有一致界。另一方面，存在 $\epsilon>0$，使每个 $|v|=2$ 的分量 $u_v$ 对能量至少贡献 $\epsilon$。具体说，考察 $0\leq i<j\leq d$ 时的所有元素对 $y_0,y_1\in\mathcal C(L_i,L_j)$，在所有为正的作用量差 $A_{H_{L_i,L_j}}(y_0)-A_{H_{L_i,L_j}}(y_1)$ 中取最小值作为 $\epsilon$。这两点合起来，就给出了\eref{9.28}中非空层所对应树 $T$ 的二价顶点数的上界。\znote{139}{$\epsilon>0$ 的存在性用到交点集的有限性：$\mathcal C(L_i,L_j)$ 是紧致 Lagrangian 子流形的横截交点集，故为有限集，于是正的作用量差只有有限多个，其最小值即为 $\epsilon$。这个下界的作用是给二价顶点数一个与树无关的上界：每插入一个二价顶点至少消耗能量 $\epsilon$，而总能量被右端的作用量差与曲率项的一致界控制，从而说明\eref{9.28}中非空层只有有限多个，$\mu^d$ 的计数才是有限和。[来源：原书 §(9l)，能量恒等式见\eref{8.12}。]}

从现在起，假设所有相关模空间都正则。用 $\M^{d+1}(y_0,\ldots,y_d)^p$ 表示 $\M^{d+1}(y_0,\ldots,y_d)$ 的 $p$ 维部分，它是该空间中的开闭子集；类似地定义 $\M^T(y_0,\ldots,y_d)^p$。\eref{9.28}上存在自然的稳定映射拓扑，在这个拓扑下整个空间紧致。并且，每个
\begin{equation}\label{eq:9.29}
\overline\M^{d+1}(y_0,\ldots,y_d)^p
=\coprod_T\M^T(y_0,\ldots,y_d)^{p-|\operatorname{Ed}^{\mathrm{int}}(T)|}
\end{equation}
都是开闭子集，同时是 $p$ 维带角拓扑流形。更精确地，单顶点树对应的部分可与 $\M^{d+1}(y_0,\ldots,y_d)^p$ 等同，它构成内部；其余部分都是余维数为 $|\operatorname{Ed}^{\mathrm{int}}(T)|$ 的边界层。\eref{9.29}的可微结构问题相当微妙：每一层当然都是光滑流形，而且可以定义产生各边界层领状邻域的粘合映射。但是，粘合参数的选择并不那么明确，因此这些映射可能不像 $\overline{\mathcal R}^{d+1}$ 的情形那样很好地拼合。幸好，构造 Fukaya 范畴只用到维数 $p\leq1$ 的模空间，此时这个问题无关紧要。

先考虑 $p=0$，可见只有单顶点树对\eref{9.29}有贡献。该空间的紧致性意味着 $\M^{d+1}(y_0,\ldots,y_d)^0$ 是有限集，从而说明定义\eref{9.23}中的点计数是合理的。再看 $p=1$，边界点来自恰有两个顶点的半稳定树，即\sr{fig:9.5}{图~9.5} 中的 $T=T(m,n)$。设
\begin{equation}\label{eq:9.30}
\begin{aligned}
(r_1,u_1,r_2,u_2)\in{}&\M^{d+2-m}(y_0,y_1,\ldots,y_n,\tilde y,y_{n+m+1},\ldots,y_d)^0\\
&\times\M^{m+1}(\tilde y,y_{n+1},\ldots,y_{n+m})^0
\end{aligned}
\end{equation}
\SourcePage{140}
是 $T(m,n)$-层中的一点；这意味着 $\tilde y\in\mathcal C(L_n,L_{n+m})$。对所有可能的 $\tilde y$ 取并后，这样的 $(r_1,u_1,r_2,u_2)$ 的个数模 $2$，恰好是表达式
\[
\mu^{d+1-m}(y_d,\ldots,y_{n+m+1},\mu^m(y_{n+m},\ldots,y_{n+1}),y_n,\ldots,y_1)
\]
中 $y_0$ 的系数。对 $m,n$ 求和便计入全部边界点，因而必得到偶数；这正是以 $\K$ 为系数的 $\Ai$ 结合律方程\eref{1.2}。
\SourceFigure{figure-9-5}{0.46730719}{9.5}
从更具体的角度考察最后这个 Gromov 紧化情形，是很有启发性的。假设有序列 $(r^i,u^i)\in\M^{d+1}(y_0,\ldots,y_d)^1$，收敛于某个边界点\eref{9.30}。这里包含两种相关但不同的现象：
\begin{itemize}
\item 当 $m=1$ 或 $m=d$ 时，是\emph{不稳定情形}。为简化记号，只讨论前者。$u^i$ 是映射 $\mathcal S_{r^i}^{d+1}\to M$，$u_1$ 是映射 $\mathcal S_{r_1}^{d+1}\to M$，其中 $r^i,r_1\in\mathcal R^{d+1}$。事实上，$r^i\to r_1$，而 $u^i$ 在紧致子集上收敛于 $u_1$。另一分量 $u_2\in\M^2(\tilde y,y_{n+1})$ 是 Floer 方程的解，它以如下方式表现为“无穷远处的气泡”。令
\[
\epsilon^i=\epsilon_{n+1}^{d+1}|_{\{r^i\}\times Z^+}:Z^+\longrightarrow\mathcal S_{r^i}^{d+1}
\]
为 $\mathcal S_{r^i}^{d+1}$ 的第 $n+1$ 个出向条带状端。则存在 $\sigma^i\to+\infty$，使重新参数化的映射 $\tilde u^i(s,t)=u^i(\epsilon^i(s+\sigma^i,t))$ 在紧致子集上收敛于 $u_2$。这个表述有意义，因为由对任意扰动数据施加的条件\eref{8.7}，各 $\tilde u^i$ 都是 Floer 方程的局部解，并且当 $i\to\infty$ 时，它们的定义域穷竭整个 $Z$。
\item 在余下的\emph{稳定情形} $1<m<d$，$r^i$ 收敛于 $\partial\overline{\mathcal R}^{d+1}$ 中的点 $(r_1,r_2)$。按拓扑的定义，这意味着存在收敛
\SourcePage{141}
序列 $r_k^i\to r_k$ 以及粘合参数序列 $\rho^i\to0$，使得对充分大的 $i$，有 $r^i=\gamma^T(\rho^i,r_1^i,r_2^i)$。将对应的带标点圆盘记为
\[
S_1^i=\mathcal S_{r_1^i}^{d+2-m},\qquad S_2^i=\mathcal S_{r_2^i}^{m+1}.
\]
于是 $S^i$ 由 $S_1^i,S_2^i$ 按\eref{8.5}粘合而成，粘合长度为 $l^i=-\log(-\rho^i)/\pi$。把 $u^i$ 限制到 $S^i$ 中由粘合继承的各部分，得到 $S_1,S_2$ 上部分定义的映射序列，它们在紧致子集上分别收敛于 $u_1,u_2$。扰动数据选择的一致性使这一表述有意义，因为它要求 $S^i$ 和 $S_k^i$ 上所用的扰动数据彼此关联。
\end{itemize}
